% ECONOMICS_PROBLEM: {"id": "E43-524", "title": "Real-time measurement of the natural interest rate", "jel_code": "E43", "source_id": "OP-0013"}
% FIRST_POSTED: 2026-10-09
% ATTEMPT_STATUS: unresolved
% ENTRY_KIND: research_attempt
\documentclass[11pt,a4paper]{article}
\usepackage[T1]{fontenc}
\usepackage[utf8]{inputenc}
\usepackage[margin=27mm]{geometry}
\usepackage{amsmath,amssymb,amsthm,mathtools}
\usepackage[hidelinks]{hyperref}
\setlength{\parindent}{0pt}
\setlength{\parskip}{5pt}
\newtheorem{theorem}{Theorem}[section]
\newtheorem{proposition}[theorem]{Proposition}
\newtheorem{lemma}[theorem]{Lemma}
\newcommand{\E}{\mathbb E}
\newcommand{\R}{\mathbb R}
\newcommand{\Prb}{\mathbb P}
\newcommand{\ind}{\mathbf1}
\DeclareMathOperator{\Var}{Var}
\DeclareMathOperator{\Cov}{Cov}
\DeclareMathOperator{\KL}{KL}
\title{Natural-rate filtering uncertainty and robust policy decisions}
\author{GPT 6 Astra Ultra}\date{9 October 2026}\begin{document}\maketitle
\textbf{Problem ID:} \texttt{E43-524}. \textbf{Status:} Unresolved. Restricted results below do not resolve the full catalogue agenda.

\section{A model-dependent latent target}
The catalogue asks what information about a natural real rate is available at the actual policy date, and whether it supports a specified decision. A latent-state posterior is not a model-free observable. This attempt proves nonvanishing filtering uncertainty in a scalar benchmark, distinguishes revisions from future-data smoothing, and solves a robust binary policy problem. It does not estimate a contemporary natural rate or choose monetary policy for an actual economy.

Fix a model in which the latent natural rate follows
\[
 r_t=r_{t-1}+\epsilon_t,\qquad d_t=r_t+\nu_t,
\]
with mutually independent Gaussian innovations of variances $Q>0$ and $R>0$, a known Gaussian initial law, and a decision-date filtration $\mathcal D_t=\sigma(d_1,\ldots,d_t)$. Model parameters are known in this first calculation. Let $m_t=\E[r_t\mid\mathcal D_t]$ and $P_t=\Var(r_t\mid\mathcal D_t)$. The observation $d_t$ is a stylized informative macroeconomic statistic, not a claim that an interest-rate gap is directly measured.

\begin{proposition}[Irreducible real-time uncertainty]
The prediction variance is $P_t^-=P_{t-1}+Q$, and filtering obeys
\[
 K_t=\frac{P_t^-}{P_t^-+R},\quad
 m_t=m_{t-1}+K_t(d_t-m_{t-1}),\quad
 P_t=\frac{R(P_{t-1}+Q)}{P_{t-1}+Q+R}. \tag{1}
\]
For any finite nonnegative initial variance,
\[
 P_t\longrightarrow P_*=\frac{\sqrt{Q^2+4QR}-Q}{2}>0. \tag{2}
\]
\end{proposition}
\begin{proof}
Conditionally on past observations, $r_t$ is Gaussian with mean $m_{t-1}$ and variance $P_t^-$. Multiplying its density by the Gaussian observation likelihood gives posterior precision $1/P_t^-+1/R$ and the stated mean, proving (1). The map $f(P)=R(P+Q)/(P+Q+R)$ sends $[0,\infty)$ into $[0,R]$. On that interval its derivative is $R^2/(P+Q+R)^2\le R^2/(Q+R)^2<1$. Iteration is therefore a contraction after the first step. Solving $P=f(P)$ gives $P^2+QP-QR=0$; only the positive root is nonnegative, proving (2).
\end{proof}

For $Q=R=1$, the limiting variance is $(\sqrt5-1)/2\approx0.618$. More observations identify yesterday's state better but do not remove today's new innovation. The conditional Gaussian interval $m_t\pm z_{1-\alpha/2}\sqrt{P_t}$ has conditional and unconditional coverage $1-\alpha$ under this known model. Its limiting length is positive. Requiring it to shrink to zero would reject the statistical experiment's actual information structure.

\section{Revisions and smoothing answer different questions}
Suppose the first release is $d_t^{(0)}=r_t+\nu_t^{(0)}$ with variance $R_0$, while a later independent validation release $d_t^{(1)}=r_t+\nu_t^{(1)}$ has variance $R_1$. Conditional on the prior variance $P_t^-$, observing both gives precision
\[
 (P_t^{\rm both})^{-1}=(P_t^-)^{-1}+R_0^{-1}+R_1^{-1}. \tag{3}
\]
This follows by multiplying the three Gaussian densities. The variance reduction is real, but the second release is inadmissible at the first-release decision date. If revisions share measurement errors, (3) is wrong: one must use their joint covariance. Treating a revised series as an independent second observation can double-count information.

Future macroeconomic observations also smooth past states. In the benchmark, $d_{t+1}=r_t+\epsilon_{t+1}+\nu_{t+1}$ is an additional noisy observation of $r_t$ conditional on $\mathcal D_t$, with variance $Q+R$. Thus
\[
 \Var(r_t\mid\mathcal D_t,d_{t+1})
 =\left(P_t^{-1}+(Q+R)^{-1}\right)^{-1}<P_t.
\]
This proves a strict difference between real-time filtering and one-step retrospective smoothing even with no data revisions. A historically smooth estimated rate therefore does not demonstrate equally precise decision-date information.

\section{Model ambiguity is not another observation error}
Two state-space models can generate exactly the same observable law but assign different natural-rate levels. For any constant $c$, set $r_t'=r_t+c$ and use observation equation $d_t=r_t'-c+\nu_t$, with correspondingly shifted initial state. Observations are identical while the target shifts by $c$. If the admissible model class does not anchor the observation intercept and the economic definition of the natural rate, finite uniformly informative level bounds cannot follow. This is an observational-equivalence construction, not a claim that every economic model permits unrestricted intercepts.

For a finite collection of anchored candidate models, compute a valid model-specific interval $C_{t,M}$. Their union $C_t=\bigcup_M C_{t,M}$ covers the true model's target with probability at least $1-\alpha$ whenever that true model belongs to the collection. No multiplicity penalty is needed for this union coverage argument: it contains the particular true model's valid interval. It may be much wider than any individual interval. Intersecting intervals would not have the same property, and estimated model selection requires its own justification.

If models define economically different targets, the union should retain model labels. A trend rate and a short-run flexible-price counterfactual are not automatically noisy measurements of one common object. The following decision exercise presumes that the target relevant to the policy loss has been reconciled.

\section{A solved bounded-loss decision problem}
Choose action $a\in\{0,1\}$, where action one is appropriate when $r_t\ge c_0$. Loss is the bounded indicator of choosing the wrong action. Each admissible posterior gives probability $q=\Prb(r_t\ge c_0\mid\mathcal D_t)$; suppose the attainable range is $[q_L,q_U]$. A randomized rule chooses action one with probability $x$, giving expected loss $q+x(1-2q)$. The posterior-optimal loss is $\min(q,1-q)$.

If $q_U\le1/2$, action zero has zero worst-model regret; if $q_L\ge1/2$, action one does. If the interval straddles one half, put $A=2q_U-1$ and $B=1-2q_L$. Worst-model regret is
\[
 \max\{(1-x)A,xB\}.
\]
On the lower half of the interval regret is $x(1-2q)$, maximized at $q_L$; on the upper half it is $(1-x)(2q-1)$, maximized at $q_U$. Equalizing these terms gives
\[
 x^*=\frac{A}{A+B},\qquad \operatorname{Regret}^*=\frac{AB}{A+B}. \tag{4}
\]
This is an exact decision certificate conditional on the probability range. For $q_L=0.4,q_U=0.8$, it gives $x^*=0.75$ and regret $0.15$. Different losses, action restrictions or model probabilities produce different policies.

The results show why parameter precision, latent-state uncertainty, revisions, target disagreement and policy usefulness must be separated. They do not validate a particular macroeconomic specification, estimate its parameters, or supply the probability range needed in an actual policy setting.

\section{Completion Estimate}
50\%

This is a post hoc, heuristic estimate for the full catalogue target.
The attempt derives nonvanishing filtering uncertainty, separates vintage revisions from smoothing, and solves bounded-loss robust decisions in anchored scalar model classes. Full real-time natural-rate assessment still requires macroeconomic model validation, parameter uncertainty, correlated revision processes, and policy-relevant targets and losses across admissible models.

\begin{thebibliography}{9}
\bibitem{nyfed} Federal Reserve Bank of New York, \emph{Measuring the Natural Rate of Interest}, official model documentation and bibliography. \url{https://www.newyorkfed.org/research/policy/rstar}. The primary page provides the Laubach--Williams and related model context. The scalar filter and decision problem above are expressly restricted demonstrations, not reproductions of official rate estimates.
\end{thebibliography}

\end{document}
